\chapter{Mở đầu}

\section{Bối cảnh và lý do chọn đề tài}
Báo cáo này tập trung vào một bài toán thị giác máy tính có ý nghĩa ứng dụng thực tế. Phần mở đầu cần nêu rõ vấn đề đặt ra, vì sao bài toán quan trọng, và tại sao nó phù hợp với phạm vi của học phần.

\begin{itemize}
    \item Mô tả ngắn gọn bài toán hoặc hiện tượng cần xử lý trong ảnh hoặc video.
    \item Nêu bối cảnh ứng dụng, người dùng mục tiêu và giá trị thực tiễn của đề tài.
    \item Giải thích vì sao đề tài cần được giải quyết bằng các kỹ thuật thị giác máy tính.
\end{itemize}

\section{Mục tiêu}
\begin{itemize}
    \item Xây dựng được một pipeline rõ ràng từ đầu vào đến đầu ra.
    \item So sánh ít nhất hai phương pháp cho một giai đoạn chính của hệ thống.
    \item Đánh giá kết quả bằng cả chỉ số định lượng và nhận xét định tính.
\end{itemize}

\section{Phạm vi và đầu ra mong đợi}
\begin{itemize}
    \item Xác định phạm vi dữ liệu, giả định đầu vào và giới hạn của hệ thống.
    \item Chỉ rõ sản phẩm cuối cùng: notebook, script thực thi, hoặc giao diện tối thiểu.
    \item Nêu các tiêu chí thành công của đề tài ở mức báo cáo học phần.
\end{itemize}

\section{Bố cục báo cáo}
\begin{itemize}
    \item Chương 2 trình bày cơ sở lý thuyết và ánh xạ kiến thức môn học với các nhiệm vụ của đề tài.
    \item Chương 3 mô tả phương pháp giải quyết và quy trình xử lý chính.
    \item Chương 4 trình bày dữ liệu, thiết lập thí nghiệm, kết quả và phân tích lỗi.
    \item Chương 5 tổng kết những gì đã đạt được, nêu hạn chế và hướng phát triển.
    \item Phụ lục cung cấp hướng dẫn chạy, cấu trúc mã nguồn và các kết quả bổ sung nếu cần.
\end{itemize}